% ─────────────────────────────────────────────────────────────────────────────
\section{Learning Reflection}
\label{sec:reflect-learning}

\subsection{Muhammad Ahmad Humayun Hanif}

Kaavish was the first project in which I was responsible for the complete system
architecture of a production-grade application---not just a module, but the
integration of an LLM pipeline, a distributed database, a WhatsApp API, an OCR
service, and a mobile frontend into a coherent whole.  The most significant skill I
gained was \textbf{system-level thinking}: understanding how a design decision in the
data model (e.g., separating Payments from Bookings) has downstream consequences for
the agent's OCR verification loop, the vendor dashboard UI, and the audit trail.

Working with LangGraph taught me that \textit{cyclic state} is fundamentally
different from linear pipeline composition---debugging a loop that can re-enter a
node after five asynchronous hops requires a completely different mental model than
debugging a sequential program.  I also developed a deep appreciation for the
Neuro-Symbolic design pattern: the security and reliability properties it provides
(no hallucinated bookings, no prompt injection) are not achievable with a pure LLM
approach.

The hardest challenge was the Roman Urdu orthographic variation problem.  I had
underestimated how much ``waqt'' versus ``tym'' versus ``vakt'' would matter for a
tokeniser trained on formal text.  Solving it through targeted data augmentation
rather than a more complex normalisation layer was a pragmatic decision I stand by,
but it highlighted the importance of domain-specific data over model complexity.

\subsection{Jazib Waqas}

Designing and implementing the backend was simultaneously the most technical and the
most frustrating part of the project.  The Optimistic Concurrency Control
implementation in Firestore was conceptually straightforward but took significant
debugging to get right---particularly understanding that Firestore transactions retry
automatically on conflict, which required careful idempotency guarantees in our
booking creation code to prevent duplicate records on retry.

I learned the value of \textbf{API contract discipline}: establishing clear input/output
schemas for each microservice (agent, OCR, booking APIs) early in the project
prevented integration failures later.  The one week we spent without a formal API
contract between the backend and frontend resulted in two weeks of alignment work
afterwards.

WhatsApp Business API integration was technically straightforward but operationally
unpredictable: Meta's review processes for webhook verification and message template
approval add real latency to development cycles.  Future projects with third-party
communication APIs should budget this time explicitly.

\subsection{Taha Hunaid Ali}

Building the React Native frontend taught me that \textbf{state management at scale}
is a qualitatively different problem from small-project UI development.  Coordinating
real-time Firestore listeners, local optimistic UI updates (showing a slot as booked
before the OCC transaction completes), and background push notification handling
required careful architecture---I rewrote the booking flow's state management twice
before arriving at a design that was both correct and maintainable.

The social features (forum, leaderboard, chat) were the most enjoyable to build
because they had immediate visual feedback and were the most enthusiastically received
by test users.  They also taught me the importance of \textbf{performance budgeting}
on mobile: unoptimised Firestore listeners can drain a phone's battery within hours,
and pagination on the leaderboard query reduced initial load time from 8 seconds to
under 2 seconds.

\subsection{Muhammad Taqi}

My primary learning from this project was about \textbf{the cost of data}.  We spent
more calendar time constructing, annotating, and augmenting the bilingual dataset
than on any other single task.  This experience permanently reframed my view of ML
projects: the model choice matters far less than the quality and coverage of training
data.  The 20-point accuracy gap between XLM-RoBERTa (which should theoretically
be better at multilingual tasks) and DistilBERT on our small custom dataset makes
this case compellingly.

Fine-tuning transformers also taught me practical lessons that textbooks gloss over:
focal loss is essential for class imbalance but its gamma parameter requires careful
tuning; evaluation should use macro-F1 (not accuracy) when classes are imbalanced;
and confidence calibration matters when the downstream system (the dialogue state
validator) relies on model confidence to decide whether to ask a clarifying question.

% ─────────────────────────────────────────────────────────────────────────────
\section{Team Reflection}
\label{sec:reflect-team}

\subsection{Collaboration and Role Distribution}

The team operated with a flat, ownership-based model: each member owned specific
modules end-to-end (backend APIs, NLU pipeline, frontend, agent), rather than
collectively owning all code.  This accelerated individual progress but created
integration risks that required dedicated ``integration weeks'' at the end of each
phase.

In hindsight, establishing a shared API contract document at the project's start---
and treating it as a living specification updated before any breaking change---would
have reduced integration friction significantly.  We introduced this practice in
Kaavish II after experiencing its absence in Kaavish I.

\subsection{Teamwork Challenges}

The most significant teamwork challenge arose when the WhatsApp Twilio sandbox number
expired two weeks before the mid-evaluation, forcing an emergency re-integration.
The team's ability to divide the problem---one member handling the new Twilio
credentials, another updating the webhook endpoint, a third preparing a fallback
demo with the web chat interface---under time pressure was a genuine test of
collaboration and emerged as one of our team's strengths.

% ─────────────────────────────────────────────────────────────────────────────
\section{Project Reflection}
\label{sec:reflect-project}

\subsection{Proposed vs.\ Achieved}

Table~\ref{tab:proposed-vs-achieved} compares the features proposed in the original
Kaavish I proposal against what was actually delivered.

\begin{table}[H]
\centering
\caption{Proposed vs.\ achieved features.}
\label{tab:proposed-vs-achieved}
\begin{tabularx}{\textwidth}{Xcc}
\toprule
\textbf{Feature} & \textbf{Proposed} & \textbf{Achieved} \\
\midrule
WhatsApp AI Receptionist (text-based)   & \cmark & \cmark \\
Bilingual NLU (English + Roman Urdu)    & \cmark & \cmark \\
OCC slot locking (no double-booking)    & \cmark & \cmark \\
OCR payment verification                & Stretch & \cmark \\
React Native mobile app                 & \cmark & \cmark \\
Vendor dashboard with analytics         & \cmark & \cmark \\
Social features (forum, chat, friends)  & \cmark & \cmark \\
Elo-based matchmaking                   & \cmark & \cmark \\
Voice-calling agent (Twilio Voice)      & Stretch & \xmark \\
Live payment gateway (JazzCash)         & Stretch & \xmark \\
Google Sheets sync                      & \cmark & Partial \\
Multi-city expansion                    & Stretch & \xmark \\
\bottomrule
\end{tabularx}
\end{table}

OCR payment verification was a stretch goal that we implemented ahead of schedule,
replacing it in the core pipeline.  Voice calling and the live payment gateway were
deprioritised in favour of polishing the text-based WhatsApp agent and the mobile
app experience---a deliberate decision based on the user survey finding that 72\% of
target users primarily communicate via WhatsApp text.

\subsection{Design Decisions Driven by Challenges}

\begin{itemize}
  \item \textbf{Neuro-Symbolic over pure LLM:} Initially planned to use a single
        LLM for both NLU and database queries.  After observing the LLM generating
        invalid Firestore queries and ``inventing'' booking confirmations in early
        prototypes, we adopted the Neuro-Symbolic architecture where the LLM only
        classifies intent and extracts entities.

  \item \textbf{DistilBERT over Qwen2.5:} Originally planned to use Qwen2.5-7B
        as the NLU backbone.  After evaluating inference latency ($>$3 seconds on
        available GPU, exceeding our WhatsApp response budget), we switched to
        fine-tuning smaller transformer models, with DistilBERT emerging as the
        winner.

  \item \textbf{Firestore over PostgreSQL:} The SRS originally listed PostgreSQL
        as the primary database.  After evaluating real-time listener requirements
        for the social layer (chat, notifications, leaderboard) and the serverless
        OCC transaction support needed for the booking pipeline, we migrated to
        Firestore.  This decision eliminated a planned data synchronisation layer.
\end{itemize}

% ─────────────────────────────────────────────────────────────────────────────
\section{Process Reflection}
\label{sec:reflect-process}

\subsection{What Worked Well}

\begin{itemize}
  \item \textbf{Phased development:} Delivering a minimal working agent (text only,
        limited intents) in December 2025 and iterating from there gave us early
        feedback and a working demo for stakeholders at every subsequent milestone.

  \item \textbf{Structured test cases:} Maintaining a public Google Sheets test case
        document throughout development (rather than writing tests post-hoc) caught
        agent regression bugs early and gave us clear evidence for the mid-evaluation
        presentation.

  \item \textbf{Real vendor engagement:} Testing the agent with actual futsal
        operators---who sent genuine booking messages in Roman Urdu---surfaced
        orthographic variations and failure modes that synthetic data generation
        did not cover.

  \item \textbf{Early SRS/SDS discipline:} Writing the SRS and SDS before coding
        prevented scope creep and gave every team member a shared vocabulary for
        system components.
\end{itemize}

\subsection{What Could Be Improved}

\begin{itemize}
  \item \textbf{API contract documentation:} As noted above, a formal shared API
        schema document from day one would have reduced integration friction.

  \item \textbf{Earlier load testing:} We discovered the Render.com free-tier
        cold-start latency issue ($\approx$8--12 seconds) only during the
        mid-evaluation rehearsal.  Load testing in January would have given us time
        to implement response caching before the demo.

  \item \textbf{Dataset annotation tooling:} Manual annotation in a shared Google
        Sheet was slow and error-prone for the Roman Urdu samples.  A proper
        annotation tool (Label Studio, Prodigy) with inter-annotator agreement
        tracking would have improved dataset quality and annotation speed.

  \item \textbf{Third-party API contingency planning:} The WhatsApp Twilio sandbox
        number expiry was avoidable with a simple API credential rotation schedule.
        Future projects should treat API credential expiry as a first-class risk with
        a documented mitigation plan.
\end{itemize}

%%% Local Variables:
%%% mode: latex
%%% TeX-master: "../report"
%%% End:
